\documentclass[10pt,a4paper,twocolumn]{article}

% -- Codificacion y fuentes --
\usepackage[utf8]{inputenc}
\usepackage[T1]{fontenc}

% -- Margenes compactos --
\usepackage[left=1cm, right=1cm, top=1.5cm, bottom=1.5cm]{geometry}

% -- Codigo fuente --
\usepackage{listings}
\usepackage{xcolor}

\definecolor{codegreen}{rgb}{0,0.5,0}
\definecolor{codegray}{rgb}{0.5,0.5,0.5}
\definecolor{codepurple}{rgb}{0.5,0,0.35}
\definecolor{backcolour}{rgb}{0.97,0.97,0.97}

\lstset{
	backgroundcolor=\color{backcolour},
	commentstyle=\color{codegreen}\itshape,
	keywordstyle=\color{blue}\bfseries,
	numberstyle=\tiny\color{codegray},
	stringstyle=\color{codepurple},
	basicstyle=\ttfamily\scriptsize,
	breakatwhitespace=false,
	breaklines=true,
	captionpos=b,
	keepspaces=true,
	numbers=left,
	numbersep=4pt,
	showspaces=false,
	showstringspaces=false,
	showtabs=false,
	tabsize=2,
	frame=single,
	rulecolor=\color{black!30},
	xleftmargin=1.5em,
	framexleftmargin=1em,
}

\lstdefinestyle{cpp}{language=C++,
	morekeywords={int64_t, uint64_t, __int128, __builtin_popcount, __gcd}
}
\lstdefinestyle{python}{language=Python}
\lstdefinestyle{java}{language=Java}

% -- Matematicas --
\usepackage{amsmath, amssymb}

% -- Tabla de contenidos clickeable --
\usepackage{hyperref}
\hypersetup{
	colorlinks=true,
	linkcolor=blue!60!black,
	urlcolor=blue!60!black,
}

% -- Headers con nombre del equipo --
\usepackage{fancyhdr}
\pagestyle{fancy}
\fancyhf{}
\fancyhead[L]{\textbf{cin>>nombre;}}
\fancyhead[C]{\textbf{ICPC Notebook}}
\fancyhead[R]{P\'agina \thepage}
\renewcommand{\headrulewidth}{0.4pt}

% -- Secciones compactas --
\usepackage{titlesec}
\titlespacing*{\section}{0pt}{1.5ex}{0.5ex}
\titlespacing*{\subsection}{0pt}{1ex}{0.3ex}
\titlespacing*{\subsubsection}{0pt}{0.8ex}{0.2ex}

\begin{document}
	
	\begin{center}
		{\LARGE \textbf{Handbook cin>>nombre;}} \\[4pt]
		{\large cin>>nombre; --- UTFSM} \\[2pt]
		{\small \'Ultima actualizaci\'on: \today}
	\end{center}
	
	\tableofcontents
	\vspace{0.5cm}
	
	% -----------------------------------------------
	\section{Template Base}
	% -----------------------------------------------
	
	\begin{lstlisting}[style=cpp, title={template.cpp}]
		#include <bits/stdc++.h>
		using namespace std;
		using ll = long long;
		#define all(v) (v).begin(), (v).end()
		
		void solve(){
			
		}
		int main(){
			ios::sync_with_stdio(false);
			cin.tie(nullptr);
			ll t = 1;
			cin >> t;
			while(t--){
				solve();
			}
			return 0;
		}
	\end{lstlisting}
	
	% -----------------------------------------------
	\section{Grafos}
	% -----------------------------------------------
	
	\subsection{BFS / DFS}
	% Tu codigo aqui
	
	\subsection{Dijkstra -- $O((V+E) \log V)$}
	
	\begin{lstlisting}[style=cpp, title={dijkstra.cpp}]
		/*
		*Description:* Shortest path, O((V+E) log V),
		no negative weights
		*/
		for (int i = 1; i <= n; i++) distance[i] = INF;
		distance[x] = 0;
		q.push({0,x});
		while (!q.empty()) {
			int a = q.top().second; q.pop();
			if (processed[a]) continue;
			processed[a] = true;
			for (auto u : adj[a]) {
				int b = u.first, w = u.second;
				if (distance[a]+w < distance[b]) {
					distance[b] = distance[a]+w;
					q.push({-distance[b],b});
				}
			}
		}
	\end{lstlisting}
	
	\subsection{Floyd-Warshall -- $O(V^3)$}
	
	\begin{lstlisting}[style=cpp, title={Floyd-Warshall.cpp}]
		/*
		*Description:* All pairs shortest path, O(V^3)
		*Status:* Tested
		*/
		for (int i = 1; i <= n; i++) {
			for (int j = 1; j <= n; j++) {
				if (i == j) distance[i][j] = 0;
				else if (adj[i][j]) distance[i][j] = adj[i][j];
				else distance[i][j] = INF;
			}
		}
		for (int k = 1; k <= n; k++) {
			for (int i = 1; i <= n; i++) {
				for (int j = 1; j <= n; j++) {
					distance[i][j] = min(distance[i][j],
					distance[i][k]+distance[k][j]);
				}
			}
		}
	\end{lstlisting}
	
	\subsection{Bellman-Ford -- $O(VE)$}
	
	\begin{lstlisting}[style=cpp, title={Bellman-Ford.cpp}]
		/*
		*Description:* Shortest paths from s,
		allows negative weights. O(VE)
		*Status:* Tested on CSES
		*/
		struct BellmanFord {
			struct Edge { int from, to; ll weight; };
			int n, last_updated = -1;
			const ll INF = 1e18;
			vector<int> p; vector<ll> dist;
			BellmanFord(vector<Edge> &G, int s, int _n) {
				n = _n; dist.assign(n+1,INF);
				p.assign(n+1,-1); dist[s] = 0;
				for (int i = 1; i <= n; i++) {
					last_updated = -1;
					for (Edge &e : G)
					if (dist[e.from]+e.weight < dist[e.to]) {
						dist[e.to] = dist[e.from]+e.weight;
						p[e.to] = e.from;
						last_updated = e.to;
					}
				}
			}
			bool getCycle(vector<int> &cycle) {
				if (last_updated == -1) return false;
				for (int i = 0; i < n-1; i++)
				last_updated = p[last_updated];
				for (int x = last_updated ;; x=p[x]) {
					cycle.push_back(x);
					if (x == last_updated
					and cycle.size() > 1) break;
				}
				reverse(cycle.begin(), cycle.end());
				return true;
			}
		};
	\end{lstlisting}
	
	\subsection{Kruskal -- $O(E \log E)$}
	
	\begin{lstlisting}[style=cpp, title={Kruskal.cpp}]
		/*
		*Description:* MST in O(E log E)
		*Status:* Tested
		*/
		struct Edge {
			int a; int b; ll w;
			Edge(int a_, int b_, ll w_)
			: a(a_), b(b_), w(w_) {}
		};
		bool c_edge(Edge &a, Edge &b) {
			return a.w < b.w;
		}
		ll Kruskal() {
			int n = G.size();
			union_find sets(n);
			vector<Edge> edges;
			for(int i = 0; i < n; i++) {
				for(auto eg : G[i]) {
					edges.emplace_back(
					i, eg.first, eg.second);
				}
			}
			sort(edges.begin(), edges.end(), c_edge);
			ll min_cost = 0;
			for(Edge e : edges) {
				if(!sets.sameSet(e.a, e.b)) {
					tree.emplace_back(e.a, e.b, e.w);
					min_cost += e.w;
					sets.unionSet(e.a, e.b);
				}
			}
			return min_cost;
		}
	\end{lstlisting}
	
	\subsection{LCA -- $O(V \log V)$ build, $O(\log V)$ query}
	
	\begin{lstlisting}[style=cpp, title={LCA.cpp}]
		/*
		*Description:* Lowest common ancestor via
		binary lifting.
		*Status:* Tested
		*/
		struct LCA {
			vector<vector<int>> T, parent;
			vector<int> depth, vis;
			int LOGN, V;
			LCA(vector<vector<int>> &T, int logn = 20)
			: T(T), LOGN(logn), vis(T.size()) {
				parent.assign(T.size()+1,
				vector<int>(LOGN, 0));
				depth.assign(T.size()+1, 0);
			}
			void dfs(int u = 0, int p = -1) {
				vis[u] = true;
				for (int v : T[u]) {
					if (p != v) {
						depth[v] = depth[u] + 1;
						parent[v][0] = u;
						for (int j = 1; j < LOGN; j++)
						parent[v][j] =
						parent[parent[v][j-1]][j-1];
						dfs(v, u);
					}
				}
			}
			int query(int u, int v) {
				if (depth[u] < depth[v]) swap(u, v);
				int k = depth[u]-depth[v];
				for (int j = LOGN-1; j >= 0; j--)
				if (k & (1 << j))
				u = parent[u][j];
				if (u == v) return u;
				for (int j = LOGN-1; j >= 0; j--) {
					if (parent[u][j] != parent[v][j]) {
						u = parent[u][j];
						v = parent[v][j];
					}
				}
				return parent[u][0];
			}
		};
	\end{lstlisting}
	
	\subsection{Union-Find / DSU -- $O(\alpha(n))$}
	
	\begin{lstlisting}[style=cpp, title={dsu.cpp}]
		/*
		*Description:* Union-Find, path compression
		+ union by size. O(alpha(n)) amortized.
		*Status:* Highly tested
		*/
		struct union_find {
			vector<int> e;
			union_find(int n) { e.assign(n, -1); }
			int findSet(int x) {
				return (e[x] < 0
				? x : e[x] = findSet(e[x]));
			}
			bool sameSet(int x, int y) {
				return findSet(x) == findSet(y);
			}
			int size(int x) { return -e[findSet(x)]; }
			bool unionSet(int x, int y) {
				x = findSet(x), y = findSet(y);
				if (x == y) return 0;
				if (e[x] > e[y]) swap(x, y);
				e[x] += e[y], e[y] = x;
				return 1;
			}
		};
	\end{lstlisting}
	
	% -----------------------------------------------
	\section{Estructuras de Datos}
	% -----------------------------------------------
	
	\subsection{Segment Tree -- $O(\log n)$}
	
	\begin{lstlisting}[style=cpp, title={segtree.cpp}]
		/*
		*Description:* Build O(n), query/update O(log n),
		positions [0, n-1]
		*Status:* Highly tested
		*/
		int sum_op(int a, int b){ return a + b; }
		
		template<class T, T (*m_)(T, T)>
		struct segment_tree{
			int n; vector<T> ST;
			segment_tree(){}
			segment_tree(vector<T> &a){
				n = a.size(); ST.resize(n << 1);
				for(int i=n;i<(n<<1);i++) ST[i]=a[i-n];
				for(int i=n-1;i>0;i--)
				ST[i]=m_(ST[i<<1],ST[i<<1|1]);
			}
			void update(int pos, T val){
				ST[pos += n] = val;
				for(pos >>= 1; pos > 0; pos >>= 1)
				ST[pos] = m_(ST[pos<<1], ST[pos<<1|1]);
			}
			T query(int l, int r){
				T ansL, ansR; bool hasL=0, hasR=0;
				for(l+=n, r+=n+1; l<r; l>>=1, r>>=1){
					if(l&1)
					ansL=(hasL?m_(ansL,ST[l++]):ST[l++]),
					hasL=1;
					if(r&1)
					ansR=(hasR?m_(ST[--r],ansR):ST[--r]),
					hasR=1;
				}
				if(!hasL) return ansR;
				if(!hasR) return ansL;
				return m_(ansL, ansR);
			}
		};
		
		int main(){
			int n; cin >> n;
			vector<int> a(n);
			for(int i = 0; i < n; i++) cin >> a[i];
			segment_tree<int, sum_op> st(a);
			int q; cin >> q;
			while(q--){
				int type; cin >> type;
				if(type == 1){
					int pos, val; cin >> pos >> val;
					st.update(pos, val);
				}else{
					int l, r; cin >> l >> r;
					cout << st.query(l, r) << '\n';
				}
			}
		}
	\end{lstlisting}
	
	\subsection{Sparse Table -- $O(n \log n)$ build, $O(1)$ query}
	
	\begin{lstlisting}[style=cpp, title={sparse-table.cpp}]
		/*
		*Description:* Build O(n log n), query O(1).
		Only idempotent ops (min, max, gcd).
		No updates.
		*Status:* Highly tested
		*/
		int min_op(int a, int b){ return min(a, b); }
		
		template <typename T, T (*m_)(T, T)>
		struct sparse_table {
			int n;
			vector<vector<T>> table;
			sparse_table() {}
			sparse_table(vector<T> &arr) {
				n = arr.size();
				int k = log2_floor(n) + 1;
				table.assign(n, vector<T>(k));
				for(int i = 0; i < n; i++)
				table[i][0] = arr[i];
				for(int j = 1; j < k; j++)
				for(int i = 0; i+(1<<j) <= n; i++)
				table[i][j] = m_(table[i][j-1],
				table[i+(1<<(j-1))][j-1]);
			}
			T query(int l, int r) {
				int k = log2_floor(r - l + 1);
				return m_(table[l][k],
				table[r-(1<<k)+1][k]);
			}
			int log2_floor(int n) {
				return n ? __builtin_clzll(1)
				- __builtin_clzll(n) : -1;
			}
		};
		
		int main(){
			int n; cin >> n;
			vector<int> a(n);
			for(int i = 0; i < n; i++) cin >> a[i];
			sparse_table<int, min_op> st(a);
			int q; cin >> q;
			while(q--){
				int l, r; cin >> l >> r;
				cout << st.query(l, r) << '\n';
			}
		}
	\end{lstlisting}
	
	\subsection{Segment Tree Lazy -- $O(\log n)$}
	
	\begin{lstlisting}[style=cpp, title={segment-tree-lazy.cpp}]
		/*
		*Description:* Seg tree + range updates,
		build O(n), query/update O(log n)
		*Status:* Highly tested
		*/
		using T1 = long long;
		using T2 = long long;
		
		T1 merge_op(T1 a, T1 b){ return a + b; }
		
		void pushUpd_op(T2& parent, T2& child,
		int l, int r, int cl, int cr){
			child += parent;
		}
		
		void applyUpd_op(T2& lazy, T1& val,
		int l, int r){
			val += (r - l + 1) * lazy;
		}
		
		template<
		class T1, class T2,
		T1 (*merge)(T1, T1),
		void (*pushUpd)(T2&,T2&,int,int,int,int),
		void (*applyUpd)(T2&, T1&, int, int)
		>
		struct segment_tree_lazy{
			vector<T1> ST; vector<T2> lazy;
			vector<bool> upd; int n;
			void build(int i, int l, int r,
			vector<T1>&values){
				if(l == r){ ST[i] = values[l]; return; }
				build(i<<1, l, (l+r)>>1, values);
				build(i<<1|1, (l+r)/2+1, r, values);
				ST[i] = merge(ST[i<<1], ST[i<<1|1]);
			}
			segment_tree_lazy() {}
			segment_tree_lazy(vector<T1>&values){
				n = values.size(); ST.resize(n<<2|3);
				lazy.resize(n<<2|3);
				upd.resize(n<<2|3, false);
				build(1, 0, n-1, values);
			}
			void push(int i, int l, int r){
				if(upd[i]){
					applyUpd(lazy[i], ST[i], l, r);
					if(l != r){
						pushUpd(lazy[i], lazy[i<<1],
						l, r, l, (l+r)/2);
						pushUpd(lazy[i], lazy[i<<1|1],
						l, r, (l+r)/2+1, r);
						upd[i<<1] = 1;
						upd[i<<1|1] = 1;
					}
					upd[i] = false;
					lazy[i] = T2();
				}
			}
			void update(int i, int l, int r,
			int a, int b, T2 &u){
				if(l >= a and r <= b){
					pushUpd(u, lazy[i], a, b, l, r);
					upd[i] = true;
				}
				push(i, l, r);
				if(l > b or r < a) return;
				if(l >= a and r <= b) return;
				update(i<<1, l, (l+r)>>1, a, b, u);
				update(i<<1|1, (l+r)/2+1, r, a, b, u);
				ST[i] = merge(ST[i<<1], ST[i<<1|1]);
			}
			void update(int a, int b, T2 u){
				if(a > b){
					update(0, b, u);
					update(a, n-1, u);
					return;
				}
				update(1, 0, n-1, a, b, u);
			}
			T1 query(int i, int l, int r,
			int a, int b){
				push(i, l, r);
				if(a <= l and r <= b) return ST[i];
				int mid = (l+r) >> 1;
				if(mid < a)
				return query(i<<1|1, mid+1, r, a, b);
				if(mid >= b)
				return query(i<<1, l, mid, a, b);
				return merge(
				query(i<<1, l, mid, a, b),
				query(i<<1|1, mid+1, r, a, b));
			}
			T1 query(int a, int b){
				if(a > b)
				return merge(query(a,n-1), query(0,b));
				return query(1, 0, n-1, a, b);
			}
		};
		
		int main(){
			int n; cin >> n;
			vector<long long> a(n);
			for(int i = 0; i < n; i++) cin >> a[i];
			segment_tree_lazy<long long, long long,
			merge_op, pushUpd_op, applyUpd_op> st(a);
			int q; cin >> q;
			while(q--){
				int type; cin >> type;
				if(type == 1){
					int l, r; long long val;
					cin >> l >> r >> val;
					st.update(l, r, val);
				}else{
					int l, r; cin >> l >> r;
					cout << st.query(l, r) << '\n';
				}
			}
		}
	\end{lstlisting}
	
	% -----------------------------------------------
	\section{Matem\'aticas}
	% -----------------------------------------------
	
	\subsection{BinPow -- $O(\log b)$}
	
	\begin{lstlisting}[style=cpp, title={BinPow.cpp}]
		/*
		*Description:* a^b mod m in O(log b)
		*Status:* Highly tested
		*/
		ll binpow(ll a, ll b, ll m) {
			a %= m;
			ll res = 1;
			while (b > 0) {
				if (b & 1)
				res = (res * a) % m;
				a = (a * a) % m;
				b >>= 1;
			}
			return res;
		}
	\end{lstlisting}
	
	\subsection{GetDivisors -- $O(\sqrt{n})$}
	
	\begin{lstlisting}[style=cpp, title={GetDivisors.cpp}]
		/*
		*Description:* All divisors in O(sqrt(n)),
		includes 1 and n
		*Status:* Highly tested
		*/
		vector<ll> get_divisors(ll n) {
			vector<ll> left, right;
			for (ll i = 1; i * i <= n; i++)
			if (n % i == 0) {
				left.push_back(i);
				if (i != n / i)
				right.push_back(n / i);
			}
			reverse(begin(right), end(right));
			for (ll x : right) left.push_back(x);
			return left;
		}
	\end{lstlisting}
	
	\subsection{Criba de Erat\'ostenes -- $O(n \log \log n)$}
	
	\begin{lstlisting}[style=cpp, title={criba.cpp}]
		/*
		*Description:* All primes up to n,
		O(n log(log(n)))
		*Status:* Highly tested
		*/
		struct eratosthenes_sieve {
			vector<ll> primes;
			vector<bool> isPrime;
			eratosthenes_sieve(ll n) {
				isPrime.resize(n + 1, true);
				isPrime[0] = isPrime[1] = false;
				for (ll i = 2; i <= n; i++) {
					if (isPrime[i]) {
						primes.push_back(i);
						for (ll j = i*i; j <= n; j += i)
						isPrime[j] = false;
					}
				}
			}
		};
	\end{lstlisting}
	
	\subsection{GCD extendido -- $O(\log(\min(a,b)))$}
	
	\begin{lstlisting}[style=cpp, title={GCD-extendido.cpp}]
		/*
		*Description:* Extended euclidean algorithm,
		O(log(min(a,b)))
		*Status:* Tested
		*/
		struct GCD_type { ll x, y, d; };
		GCD_type ex_GCD(ll a, ll b){
			if (b == 0) return {1, 0, a};
			GCD_type pom = ex_GCD(b, a % b);
			return {pom.y, pom.x - a/b * pom.y, pom.d};
		}
	\end{lstlisting}
	
	\subsection{Factores Primos -- $O(\sqrt{n})$}
	
	\begin{lstlisting}[style=cpp, title={Factores-primos.cpp}]
		/*
		*Description:* Prime factors of n, O(sqrt(n))
		*Status:* Highly tested
		*/
		vector<ll> primeFactors(ll n) {
			vector<ll> factors;
			for (ll i = 2; (i*i) <= n; i++) {
				while (n % i == 0) {
					factors.push_back(i);
					n /= i;
				}
			}
			if (n > 1) factors.push_back(n);
			return factors;
		}
	\end{lstlisting}
	
	% -----------------------------------------------
	\section{Strings}
	% -----------------------------------------------
	
	\subsection{KMP -- $O(n+m)$}
	
	\begin{lstlisting}[style=cpp, title={KMP.cpp}]
		/*
		*Description:* Pattern matching in O(n+m)
		*Status:* Tested on CSES
		*/
		template<class T> struct KMP {
			T pattern; vector<int> lps;
			KMP(T &pat): pattern(pat) {
				lps.resize(pat.size(), 0);
				int len = 0, i = 1;
				while (i < pattern.size()) {
					if (pattern[i] == pattern[len])
					lps[i++] = ++len;
					else {
						if (len != 0) len = lps[len-1];
						else lps[i++] = 0;
					}
				}
			}
			vector<int> search(T &text) {
				vector<int> matches;
				int i = 0, j = 0;
				while (i < text.size()) {
					if (pattern[j] == text[i]) {
						i++, j++;
						if (j == pattern.size()) {
							matches.push_back(i - j);
							j = lps[j - 1];
						}
					} else {
						if (j != 0) j = lps[j-1];
						else i++;
					}
				}
				return matches;
			}
		};
		
		// Prefix function (simple version)
		template<class T>
		vector<int> prefix(T &S) {
			vector<int> P(S.size());
			P[0] = 0;
			for(int i = 1; i < S.size(); ++i) {
				P[i] = P[i - 1];
				while(P[i] > 0 && S[P[i]] != S[i])
				P[i] = P[P[i] - 1];
				if(S[P[i]] == S[i])
				++P[i];
			}
			return P;
		}
		
		int main(){
			string text, pattern;
			cin >> text >> pattern;
			KMP<string> kmp(pattern);
			vector<int> matches = kmp.search(text);
			cout << matches.size() << '\n';
			for(int x : matches) cout << x << " ";
			cout << '\n';
		}
	\end{lstlisting}
	
	\subsection{Z-Function -- $O(n+m)$}
	
	\begin{lstlisting}[style=cpp, title={z-function.cpp}]
		/*
		*Description:* Z-algorithm, all occurrences
		in O(n+m)
		*Status:* Tested
		*/
		struct Z {
			int n, m;
			vector<int> z;
			Z(string s) {
				n = s.size();
				z.assign(n, 0);
				int l = 0, r = 0;
				for (int i = 1; i < n; i++) {
					if (i <= r)
					z[i] = min(r-i+1, z[i-l]);
					while (i+z[i] < n
					&& s[z[i]] == s[i+z[i]])
					++z[i];
					if (i+z[i]-1 > r)
					l = i, r = i+z[i]-1;
				}
			}
			Z(string p, string t) {
				string s = p + "#" + t;
				n = p.size();
				m = t.size();
				z.assign(n+m+1, 0);
				int l = 0, r = 0;
				for (int i = 1; i < n+m+1; i++) {
					if (i <= r)
					z[i] = min(r-i+1, z[i-l]);
					while (i+z[i] < n+m+1
					&& s[z[i]] == s[i+z[i]])
					++z[i];
					if (i+z[i]-1 > r)
					l = i, r = i+z[i]-1;
				}
			}
			void p_in_t(vector<int>& ans) {
				for (int i = n+1; i < n+m+1; i++)
				if (z[i] == n)
				ans.push_back(i-n-1);
			}
		};
		
		int main(){
			string t, p;
			cin >> t >> p;
			Z z(p, t);
			vector<int> ans;
			z.p_in_t(ans);
			cout << ans.size() << '\n';
			for(int x : ans) cout << x << " ";
			cout << '\n';
		}
	\end{lstlisting}
	
	\subsection{Trie -- $O(|s|)$ por operaci\'on}
	
	\begin{lstlisting}[style=cpp, title={trie.cpp}]
		/*
		*Description:* Trie for prefix queries.
		O(|s|) per operation.
		*Status:* Tested
		*/
		struct trie {
			vector<vector<ll>> tree;
			vector<ll> counts;
			trie() {
				tree.push_back(vector<ll>(26, -1));
				counts.push_back(0);
			};
			void insert(string &s) {
				ll i = 0, u = 0;
				while(i < s.size()){
					char c = s[i];
					if (tree[u][c-'a'] == -1){
						ll pos = tree.size();
						tree.push_back(vector<ll>(26, -1));
						counts.push_back(0);
						tree[u][c-'a'] = pos;
					}
					i++;
					u = tree[u][c-'a'];
				}
				counts[u]++;
			}
			ll search(string &s){
				ll i = 0, u = 0;
				ll count_ends = 0;
				while(i < s.size()){
					char c = s[i];
					if (tree[u][c-'a'] != -1){
						u = tree[u][c-'a'];
						i++;
					}else break;
				}
				if(i == s.size())
				count_ends += counts[u];
				return count_ends;
			}
		};
		
		int main(){
			int n; cin >> n;
			trie t;
			for(int i = 0; i < n; i++){
				string s; cin >> s;
				t.insert(s);
			}
			int q; cin >> q;
			while(q--){
				string s; cin >> s;
				cout << t.search(s) << '\n';
			}
		}
	\end{lstlisting}
	
	% -----------------------------------------------
	\section{Geometr\'ia}
	% -----------------------------------------------
	% Tu codigo aqui
	
	% -----------------------------------------------
	\section{Programaci\'on Din\'amica}
	% -----------------------------------------------
	
	\subsection{Longest Increasing Subsequence -- $O(n \log n)$}
	
	\begin{lstlisting}[style=cpp, title={LIS.cpp}]
		/*
		*Description:* LIS in O(n log n),
		allows sequence recovery
		*Status:* Highly tested
		*/
		template <class I>
		vector<int> LIS(const vector<I> &S) {
			if (S.empty()) return {};
			vector<int> prev(S.size());
			vector<pair<I, int>> res;
			for (int i = 0; i < S.size(); i++) {
				auto it = lower_bound(res.begin(),
				res.end(), pair<I,int>{S[i], i});
				if (it == res.end())
				res.emplace_back(), it = res.end()-1;
				*it = {S[i], i};
				prev[i] = (it == res.begin()
				? 0 : (it-1)->second);
			}
			int L = res.size(),
			cur = res.back().second;
			vector<int> ans(L);
			while (L--)
			ans[L] = cur, cur = prev[cur];
			return ans;
		}
		
		int main(){
			int n; cin >> n;
			vector<int> a(n);
			for(int i = 0; i < n; i++) cin >> a[i];
			vector<int> lis_idx = LIS(a);
			cout << lis_idx.size() << '\n';
			for(int i = 0; i < (int)lis_idx.size(); i++)
			cout << lis_idx[i] << " ";
			cout << '\n';
		}
	\end{lstlisting}
	
	\subsection{Knapsack 0/1 -- $O(nW)$}
	
	\begin{lstlisting}[style=cpp, title={mochila.cpp}]
		void solve(){
			ll n, w;
			cin >> n >> w;
			vector<pair<ll,ll>> v(n);
			for(auto &[a,b] : v) cin >> a >> b;
			vector<ll> ans(w+2,0);
			ll ans_ = 0;
			for(auto [vi, wi] : v){
				for(ll j = w; j >= wi; j--){
					ans[j] = max(ans[j], ans[j-wi]+vi);
					ans_ = max(ans_, ans[j]);
				}
			}
			cout << ans_ << endl;
		}
	\end{lstlisting}
	
	% -----------------------------------------------
	\section{Miscel\'aneos}
	% -----------------------------------------------
	
	\subsection{B\'usqueda Binaria -- $O(\log n)$}
	
	\begin{lstlisting}[style=cpp, title={binary search.cpp}]
		int bs(vector<ll>& v, int b){
			int l = 0, r = v.size();
			while(l < r){
				int mid = (l + r) / 2;
				if(v[mid] >= b) r = mid;
				else l = mid + 1;
			}
			return l;
		}
	\end{lstlisting}
	
\end{document}